\documentclass[12pt,letterpaper]{article}
\usepackage[margin=0.65in]{geometry}
\usepackage[T1]{fontenc}
\usepackage{lmodern}
\pdfmapfile{+lm.map}
\usepackage{tikz}
\pagestyle{empty}
\setlength{\parindent}{0pt}
\hyphenpenalty=10000\exhyphenpenalty=10000
\begin{document}
\begin{tikzpicture}[remember picture,overlay]\begin{scope}[shift={(current page.north west)},x=1in,y=-1in]
\node[anchor=north west,inner sep=0pt,text width=7.2in,align=left,font=\fontsize{11}{13.75}\selectfont] at (0.65,0.41) {Week 19 / No card inside another / Grades 4--5};
\draw[line width=.35pt] (.65,.70)--(7.85,.70);
\node[anchor=north west,inner sep=0pt,text width=6.7in,align=left,font=\fontsize{9}{11.25}\selectfont] at (0.65,10.43) {Bellingham Math Circle / Week 19 / F19-45-v2};
\node[anchor=north west,inner sep=0pt,text width=0.55in,align=right,font=\fontsize{9}{11.25}\selectfont] at (7.3,10.43) {1};
\node[anchor=north west,inner sep=0pt,text width=7.2in,align=left,font=\fontsize{13}{16.25}\selectfont] at (0.65,0.96) {A card fits inside another when all its pictures are on the other card. The empty card fits inside every card. In a collection, no card may fit inside another. Use a card only once in a collection.};
\node[anchor=north west,inner sep=0pt,text width=7.2in,align=left,font=\fontsize{14}{17.5}\selectfont] at (0.65,2.14) {\textbf{Problem 1:} For each deck, find the largest collection you can. Explain why no larger collection can work.};
\draw[line width=.7pt,rounded corners=2pt] (2.1550000000000002,3.1) rectangle (2.9850000000000003,3.93);
\draw[line width=1.5pt] (2.2214,3.1) -- (2.2878000000000003,3.1);
\draw[line width=.7pt,rounded corners=2pt] (3.2750000000000004,3.1) rectangle (4.105,3.93);
\draw[line width=1.5pt] (3.3414,3.1) -- (3.4078000000000004,3.1);
\draw[line width=.95pt] (3.5074000000000005,3.3324000000000003) circle[radius=0.10125999999999999in];
\draw[line width=.95pt,line join=round] (3.8726000000000003,3.2159510000000004)--(3.9839860000000002,3.4184710000000003)--(3.7612140000000003,3.4184710000000003)--cycle;
\draw[line width=.7pt,rounded corners=2pt] (4.395,3.1) rectangle (5.225,3.93);
\draw[line width=1.5pt] (4.461399999999999,3.1) -- (4.527799999999999,3.1);
\draw[line width=.95pt] (4.6274,3.3324000000000003) circle[radius=0.10125999999999999in];
\draw[line width=.7pt,rounded corners=2pt] (5.515,3.1) rectangle (6.345,3.93);
\draw[line width=1.5pt] (5.5813999999999995,3.1) -- (5.647799999999999,3.1);
\draw[line width=.95pt,line join=round] (6.1126,3.2159510000000004)--(6.223986,3.4184710000000003)--(6.001213999999999,3.4184710000000003)--cycle;
\draw[black!22,line width=.35pt] (.65,5.57)--(7.85,5.57);
\draw[line width=.7pt,rounded corners=2pt] (0.79,5.89) rectangle (1.48,6.58);
\draw[line width=1.5pt] (0.8452000000000001,5.89) -- (0.9004000000000001,5.89);
\draw[line width=.95pt] (0.9832000000000001,6.0832) circle[radius=0.08417999999999999in];
\draw[line width=.95pt] (0.89902,6.30262) rectangle (1.06738,6.47098);
\draw[line width=.7pt,rounded corners=2pt] (1.68,5.89) rectangle (2.37,6.58);
\draw[line width=1.5pt] (1.7351999999999999,5.89) -- (1.7904,5.89);
\draw[line width=.7pt,rounded corners=2pt] (2.57,5.89) rectangle (3.26,6.58);
\draw[line width=1.5pt] (2.6252,5.89) -- (2.6803999999999997,5.89);
\draw[line width=.95pt,line join=round] (3.0667999999999997,5.986393)--(3.159398,6.1547529999999995)--(2.9742019999999996,6.1547529999999995)--cycle;
\draw[line width=.7pt,rounded corners=2pt] (3.46,5.89) rectangle (4.15,6.58);
\draw[line width=1.5pt] (3.5152,5.89) -- (3.5704,5.89);
\draw[line width=.95pt] (3.6532,6.0832) circle[radius=0.08417999999999999in];
\draw[line width=.95pt,line join=round] (3.9568,5.986393)--(4.049398,6.1547529999999995)--(3.8642019999999997,6.1547529999999995)--cycle;
\draw[line width=.95pt] (3.56902,6.30262) rectangle (3.73738,6.47098);
\draw[line width=.7pt,rounded corners=2pt] (4.35,5.89) rectangle (5.039999999999999,6.58);
\draw[line width=1.5pt] (4.4052,5.89) -- (4.4604,5.89);
\draw[line width=.95pt] (4.5432,6.0832) circle[radius=0.08417999999999999in];
\draw[line width=.7pt,rounded corners=2pt] (5.239999999999999,5.89) rectangle (5.93,6.58);
\draw[line width=1.5pt] (5.2951999999999995,5.89) -- (5.3504,5.89);
\draw[line width=.95pt,line join=round] (5.7368,5.986393)--(5.829397999999999,6.1547529999999995)--(5.644202,6.1547529999999995)--cycle;
\draw[line width=.95pt] (5.349019999999999,6.30262) rectangle (5.517379999999999,6.47098);
\draw[line width=.7pt,rounded corners=2pt] (6.13,5.89) rectangle (6.82,6.58);
\draw[line width=1.5pt] (6.1852,5.89) -- (6.2404,5.89);
\draw[line width=.95pt] (6.3232,6.0832) circle[radius=0.08417999999999999in];
\draw[line width=.95pt,line join=round] (6.6268,5.986393)--(6.719398,6.1547529999999995)--(6.5342020000000005,6.1547529999999995)--cycle;
\draw[line width=.7pt,rounded corners=2pt] (7.02,5.89) rectangle (7.709999999999999,6.58);
\draw[line width=1.5pt] (7.0752,5.89) -- (7.1304,5.89);
\draw[line width=.95pt] (7.12902,6.30262) rectangle (7.2973799999999995,6.47098);
\end{scope}\end{tikzpicture}\mbox{}
\newpage
\begin{tikzpicture}[remember picture,overlay]\begin{scope}[shift={(current page.north west)},x=1in,y=-1in]
\node[anchor=north west,inner sep=0pt,text width=7.2in,align=left,font=\fontsize{11}{13.75}\selectfont] at (0.65,0.41) {Week 19 / No card inside another / Grades 4--5};
\draw[line width=.35pt] (.65,.70)--(7.85,.70);
\node[anchor=north west,inner sep=0pt,text width=6.7in,align=left,font=\fontsize{9}{11.25}\selectfont] at (0.65,10.43) {Bellingham Math Circle / Week 19 / F19-45-v2};
\node[anchor=north west,inner sep=0pt,text width=0.55in,align=right,font=\fontsize{9}{11.25}\selectfont] at (7.3,10.43) {2};
\node[anchor=north west,inner sep=0pt,text width=7.2in,align=left,font=\fontsize{14}{17.5}\selectfont] at (0.65,0.98) {\textbf{Problem 2:} Use this sixteen-card deck. Find the largest collection you can, and decide whether you can rule out every larger collection.};
\draw[line width=.7pt,rounded corners=2pt] (0.79,2.0) rectangle (1.48,2.69);
\draw[line width=1.5pt] (0.8452000000000001,2.0) -- (0.9004000000000001,2.0);
\draw[line width=.95pt] (0.9832000000000001,2.1932) circle[radius=0.08417999999999999in];
\draw[line width=.95pt] (0.89902,2.41262) rectangle (1.06738,2.58098);
\draw[line width=.7pt,rounded corners=2pt] (1.68,2.0) rectangle (2.37,2.69);
\draw[line width=1.5pt] (1.7351999999999999,2.0) -- (1.7904,2.0);
\draw[line width=.95pt,line join=round] (2.1768,2.096393)--(2.2693980000000002,2.2647530000000002)--(2.084202,2.2647530000000002)--cycle;
\draw[line width=.85pt,line join=round] (2.17680,2.38737)--(2.20352,2.46002)--(2.28088,2.46298)--(2.22003,2.51085)--(2.24112,2.58533)--(2.17680,2.54226)--(2.11248,2.58533)--(2.13357,2.51085)--(2.07272,2.46298)--(2.15008,2.46002)--cycle;
\draw[line width=.7pt,rounded corners=2pt] (2.57,2.0) rectangle (3.26,2.69);
\draw[line width=1.5pt] (2.6252,2.0) -- (2.6803999999999997,2.0);
\draw[line width=.7pt,rounded corners=2pt] (3.46,2.0) rectangle (4.15,2.69);
\draw[line width=1.5pt] (3.5152,2.0) -- (3.5704,2.0);
\draw[line width=.95pt] (3.6532,2.1932) circle[radius=0.08417999999999999in];
\draw[line width=.95pt,line join=round] (3.9568,2.096393)--(4.049398,2.2647530000000002)--(3.8642019999999997,2.2647530000000002)--cycle;
\draw[line width=.95pt] (3.56902,2.41262) rectangle (3.73738,2.58098);
\draw[line width=.7pt,rounded corners=2pt] (4.35,2.0) rectangle (5.039999999999999,2.69);
\draw[line width=1.5pt] (4.4052,2.0) -- (4.4604,2.0);
\draw[line width=.95pt] (4.45902,2.41262) rectangle (4.62738,2.58098);
\draw[line width=.85pt,line join=round] (4.84680,2.38737)--(4.87352,2.46002)--(4.95088,2.46298)--(4.89003,2.51085)--(4.91112,2.58533)--(4.84680,2.54226)--(4.78248,2.58533)--(4.80357,2.51085)--(4.74272,2.46298)--(4.82008,2.46002)--cycle;
\draw[line width=.7pt,rounded corners=2pt] (5.239999999999999,2.0) rectangle (5.93,2.69);
\draw[line width=1.5pt] (5.2951999999999995,2.0) -- (5.3504,2.0);
\draw[line width=.95pt] (5.433199999999999,2.1932) circle[radius=0.08417999999999999in];
\draw[line width=.7pt,rounded corners=2pt] (6.13,2.0) rectangle (6.82,2.69);
\draw[line width=1.5pt] (6.1852,2.0) -- (6.2404,2.0);
\draw[line width=.95pt,line join=round] (6.6268,2.096393)--(6.719398,2.2647530000000002)--(6.5342020000000005,2.2647530000000002)--cycle;
\draw[line width=.95pt] (6.23902,2.41262) rectangle (6.40738,2.58098);
\draw[line width=.85pt,line join=round] (6.62680,2.38737)--(6.65352,2.46002)--(6.73088,2.46298)--(6.67003,2.51085)--(6.69112,2.58533)--(6.62680,2.54226)--(6.56248,2.58533)--(6.58357,2.51085)--(6.52272,2.46298)--(6.60008,2.46002)--cycle;
\draw[line width=.7pt,rounded corners=2pt] (7.02,2.0) rectangle (7.709999999999999,2.69);
\draw[line width=1.5pt] (7.0752,2.0) -- (7.1304,2.0);
\draw[line width=.95pt] (7.2132,2.1932) circle[radius=0.08417999999999999in];
\draw[line width=.95pt,line join=round] (7.5168,2.096393)--(7.609398,2.2647530000000002)--(7.424202,2.2647530000000002)--cycle;
\draw[line width=.7pt,rounded corners=2pt] (0.79,2.8899999999999997) rectangle (1.48,3.5799999999999996);
\draw[line width=1.5pt] (0.8452000000000001,2.8899999999999997) -- (0.9004000000000001,2.8899999999999997);
\draw[line width=.85pt,line join=round] (1.28680,3.27737)--(1.31352,3.35002)--(1.39088,3.35298)--(1.33003,3.40085)--(1.35112,3.47533)--(1.28680,3.43226)--(1.22248,3.47533)--(1.24357,3.40085)--(1.18272,3.35298)--(1.26008,3.35002)--cycle;
\draw[line width=.7pt,rounded corners=2pt] (1.68,2.8899999999999997) rectangle (2.37,3.5799999999999996);
\draw[line width=1.5pt] (1.7351999999999999,2.8899999999999997) -- (1.7904,2.8899999999999997);
\draw[line width=.95pt,line join=round] (2.1768,2.9863929999999996)--(2.2693980000000002,3.154753)--(2.084202,3.154753)--cycle;
\draw[line width=.95pt] (1.78902,3.3026199999999997) rectangle (1.95738,3.4709799999999995);
\draw[line width=.7pt,rounded corners=2pt] (2.57,2.8899999999999997) rectangle (3.26,3.5799999999999996);
\draw[line width=1.5pt] (2.6252,2.8899999999999997) -- (2.6803999999999997,2.8899999999999997);
\draw[line width=.95pt] (2.7632,3.0831999999999997) circle[radius=0.08417999999999999in];
\draw[line width=.95pt,line join=round] (3.0667999999999997,2.9863929999999996)--(3.159398,3.154753)--(2.9742019999999996,3.154753)--cycle;
\draw[line width=.95pt] (2.67902,3.3026199999999997) rectangle (2.84738,3.4709799999999995);
\draw[line width=.85pt,line join=round] (3.06680,3.27737)--(3.09352,3.35002)--(3.17088,3.35298)--(3.11003,3.40085)--(3.13112,3.47533)--(3.06680,3.43226)--(3.00248,3.47533)--(3.02357,3.40085)--(2.96272,3.35298)--(3.04008,3.35002)--cycle;
\draw[line width=.7pt,rounded corners=2pt] (3.46,2.8899999999999997) rectangle (4.15,3.5799999999999996);
\draw[line width=1.5pt] (3.5152,2.8899999999999997) -- (3.5704,2.8899999999999997);
\draw[line width=.95pt] (3.6532,3.0831999999999997) circle[radius=0.08417999999999999in];
\draw[line width=.85pt,line join=round] (3.95680,3.27737)--(3.98352,3.35002)--(4.06088,3.35298)--(4.00003,3.40085)--(4.02112,3.47533)--(3.95680,3.43226)--(3.89248,3.47533)--(3.91357,3.40085)--(3.85272,3.35298)--(3.93008,3.35002)--cycle;
\draw[line width=.7pt,rounded corners=2pt] (4.35,2.8899999999999997) rectangle (5.039999999999999,3.5799999999999996);
\draw[line width=1.5pt] (4.4052,2.8899999999999997) -- (4.4604,2.8899999999999997);
\draw[line width=.95pt,line join=round] (4.8468,2.9863929999999996)--(4.939398,3.154753)--(4.754202,3.154753)--cycle;
\draw[line width=.7pt,rounded corners=2pt] (5.239999999999999,2.8899999999999997) rectangle (5.93,3.5799999999999996);
\draw[line width=1.5pt] (5.2951999999999995,2.8899999999999997) -- (5.3504,2.8899999999999997);
\draw[line width=.95pt] (5.433199999999999,3.0831999999999997) circle[radius=0.08417999999999999in];
\draw[line width=.95pt] (5.349019999999999,3.3026199999999997) rectangle (5.517379999999999,3.4709799999999995);
\draw[line width=.85pt,line join=round] (5.73680,3.27737)--(5.76352,3.35002)--(5.84088,3.35298)--(5.78003,3.40085)--(5.80112,3.47533)--(5.73680,3.43226)--(5.67248,3.47533)--(5.69357,3.40085)--(5.63272,3.35298)--(5.71008,3.35002)--cycle;
\draw[line width=.7pt,rounded corners=2pt] (6.13,2.8899999999999997) rectangle (6.82,3.5799999999999996);
\draw[line width=1.5pt] (6.1852,2.8899999999999997) -- (6.2404,2.8899999999999997);
\draw[line width=.95pt] (6.23902,3.3026199999999997) rectangle (6.40738,3.4709799999999995);
\draw[line width=.7pt,rounded corners=2pt] (7.02,2.8899999999999997) rectangle (7.709999999999999,3.5799999999999996);
\draw[line width=1.5pt] (7.0752,2.8899999999999997) -- (7.1304,2.8899999999999997);
\draw[line width=.95pt] (7.2132,3.0831999999999997) circle[radius=0.08417999999999999in];
\draw[line width=.95pt,line join=round] (7.5168,2.9863929999999996)--(7.609398,3.154753)--(7.424202,3.154753)--cycle;
\draw[line width=.85pt,line join=round] (7.51680,3.27737)--(7.54352,3.35002)--(7.62088,3.35298)--(7.56003,3.40085)--(7.58112,3.47533)--(7.51680,3.43226)--(7.45248,3.47533)--(7.47357,3.40085)--(7.41272,3.35298)--(7.49008,3.35002)--cycle;
\end{scope}\end{tikzpicture}\mbox{}
\newpage
\begin{tikzpicture}[remember picture,overlay]\begin{scope}[shift={(current page.north west)},x=1in,y=-1in]
\node[anchor=north west,inner sep=0pt,text width=7.2in,align=left,font=\fontsize{11}{13.75}\selectfont] at (0.65,0.41) {Week 19 / No card inside another / Grades 4--5};
\draw[line width=.35pt] (.65,.70)--(7.85,.70);
\node[anchor=north west,inner sep=0pt,text width=6.7in,align=left,font=\fontsize{9}{11.25}\selectfont] at (0.65,10.43) {Bellingham Math Circle / Week 19 / F19-45-v2};
\node[anchor=north west,inner sep=0pt,text width=0.55in,align=right,font=\fontsize{9}{11.25}\selectfont] at (7.3,10.43) {3};
\node[anchor=north west,inner sep=0pt,text width=7.2in,align=left,font=\fontsize{14}{17.5}\selectfont] at (0.65,0.98) {\textbf{Problem 3:} Use the sixteen-card deck. Keep each starting collection and add as many cards as possible. Explain whether a collection that cannot take another card must be as large as possible.};
\draw[line width=.7pt,rounded corners=2pt] (0.97,2.28) rectangle (1.88,3.19);
\draw[line width=1.5pt] (1.0428,2.28) -- (1.1156,2.28);
\draw[line width=.95pt] (1.2248,2.5347999999999997) circle[radius=0.11102000000000001in];
\draw[line width=.95pt,line join=round] (1.6252,2.4071269999999996)--(1.747322,2.629167)--(1.503078,2.629167)--cycle;
\draw[line width=.95pt] (1.1137800000000002,2.82418) rectangle (1.33582,3.04622);
\draw[line width=.85pt,line join=round] (1.62520,2.79087)--(1.66044,2.88670)--(1.76246,2.89060)--(1.68222,2.95373)--(1.71003,3.05196)--(1.62520,2.99515)--(1.54037,3.05196)--(1.56818,2.95373)--(1.48794,2.89060)--(1.58996,2.88670)--cycle;
\draw[line width=.7pt,rounded corners=2pt] (0.97,4.8) rectangle (1.88,5.71);
\draw[line width=1.5pt] (1.0428,4.8) -- (1.1156,4.8);
\draw[line width=.95pt] (1.2248,5.0548) circle[radius=0.11102000000000001in];
\draw[line width=.7pt,rounded corners=2pt] (2.16,4.8) rectangle (3.0700000000000003,5.71);
\draw[line width=1.5pt] (2.2328,4.8) -- (2.3056,4.8);
\draw[line width=.95pt,line join=round] (2.8152,4.9271270000000005)--(2.937322,5.149167)--(2.693078,5.149167)--cycle;
\draw[line width=.95pt] (2.30378,5.34418) rectangle (2.52582,5.5662199999999995);
\draw[line width=.85pt,line join=round] (2.81520,5.31087)--(2.85044,5.40670)--(2.95246,5.41060)--(2.87222,5.47373)--(2.90003,5.57196)--(2.81520,5.51515)--(2.73037,5.57196)--(2.75818,5.47373)--(2.67794,5.41060)--(2.77996,5.40670)--cycle;
\draw[line width=.7pt,rounded corners=2pt] (0.97,7.33) rectangle (1.88,8.24);
\draw[line width=1.5pt] (1.0428,7.33) -- (1.1156,7.33);
\draw[line width=.95pt] (1.2248,7.5848) circle[radius=0.11102000000000001in];
\draw[line width=.95pt,line join=round] (1.6252,7.457127000000001)--(1.747322,7.6791670000000005)--(1.503078,7.6791670000000005)--cycle;
\draw[line width=.7pt,rounded corners=2pt] (2.16,7.33) rectangle (3.0700000000000003,8.24);
\draw[line width=1.5pt] (2.2328,7.33) -- (2.3056,7.33);
\draw[line width=.95pt] (2.4148,7.5848) circle[radius=0.11102000000000001in];
\draw[line width=.95pt] (2.30378,7.87418) rectangle (2.52582,8.09622);
\end{scope}\end{tikzpicture}\mbox{}
\newpage
\begin{tikzpicture}[remember picture,overlay]\begin{scope}[shift={(current page.north west)},x=1in,y=-1in]
\node[anchor=north west,inner sep=0pt,text width=7.2in,align=left,font=\fontsize{11}{13.75}\selectfont] at (0.65,0.41) {Week 19 / No card inside another / Grades 4--5};
\draw[line width=.35pt] (.65,.70)--(7.85,.70);
\node[anchor=north west,inner sep=0pt,text width=6.7in,align=left,font=\fontsize{9}{11.25}\selectfont] at (0.65,10.43) {Bellingham Math Circle / Week 19 / F19-45-v2};
\node[anchor=north west,inner sep=0pt,text width=0.55in,align=right,font=\fontsize{9}{11.25}\selectfont] at (7.3,10.43) {4};
\node[anchor=north west,inner sep=0pt,text width=7.2in,align=left,font=\fontsize{14}{17.5}\selectfont] at (0.65,0.98) {\textbf{Problem 4:} For each deck, arrange every card into exactly one row, with each card fitting inside the next card in its row. A card can be alone in a row. Make as few rows as possible.};
\draw[line width=.7pt,rounded corners=2pt] (1.0,2.06) rectangle (1.62,2.68);
\draw[line width=1.5pt] (1.0496,2.06) -- (1.0992,2.06);
\draw[line width=.95pt] (1.1736,2.2336) circle[radius=0.07564in];
\draw[line width=.7pt,rounded corners=2pt] (2.43,2.06) rectangle (3.0500000000000003,2.68);
\draw[line width=1.5pt] (2.4796,2.06) -- (2.5292000000000003,2.06);
\draw[line width=.95pt] (2.6036,2.2336) circle[radius=0.07564in];
\draw[line width=.95pt,line join=round] (2.8764000000000003,2.146614)--(2.959604,2.297894)--(2.7931960000000005,2.297894)--cycle;
\draw[-stealth,line width=.8pt] (1.74,2.37)--(2.30,2.37);
\node[anchor=north west,inner sep=0pt,text width=4.1in,align=left,font=\fontsize{12}{15.0}\selectfont] at (3.4,2.02) {Example of one nested row: the first card fits inside the next. This is not a whole-deck arrangement.};
\draw[line width=.7pt,rounded corners=2pt] (2.25,3.17) rectangle (3.04,3.96);
\draw[line width=1.5pt] (2.3132,3.17) -- (2.3764,3.17);
\draw[line width=.7pt,rounded corners=2pt] (3.3200000000000003,3.17) rectangle (4.11,3.96);
\draw[line width=1.5pt] (3.3832000000000004,3.17) -- (3.4464,3.17);
\draw[line width=.95pt] (3.5412000000000003,3.3912) circle[radius=0.09638000000000001in];
\draw[line width=.95pt,line join=round] (3.8888000000000003,3.280363)--(3.9948180000000004,3.473123)--(3.782782,3.473123)--cycle;
\draw[line width=.7pt,rounded corners=2pt] (4.390000000000001,3.17) rectangle (5.180000000000001,3.96);
\draw[line width=1.5pt] (4.453200000000001,3.17) -- (4.516400000000001,3.17);
\draw[line width=.95pt] (4.6112,3.3912) circle[radius=0.09638000000000001in];
\draw[line width=.7pt,rounded corners=2pt] (5.46,3.17) rectangle (6.25,3.96);
\draw[line width=1.5pt] (5.5232,3.17) -- (5.5864,3.17);
\draw[line width=.95pt,line join=round] (6.0288,3.280363)--(6.134818,3.473123)--(5.922782000000001,3.473123)--cycle;
\draw[black!22,line width=.35pt] (.65,4.98)--(7.85,4.98);
\draw[line width=.7pt,rounded corners=2pt] (0.79,5.29) rectangle (1.48,5.98);
\draw[line width=1.5pt] (0.8452000000000001,5.29) -- (0.9004000000000001,5.29);
\draw[line width=.95pt] (0.9832000000000001,5.4832) circle[radius=0.08417999999999999in];
\draw[line width=.95pt] (0.89902,5.7026200000000005) rectangle (1.06738,5.87098);
\draw[line width=.7pt,rounded corners=2pt] (1.68,5.29) rectangle (2.37,5.98);
\draw[line width=1.5pt] (1.7351999999999999,5.29) -- (1.7904,5.29);
\draw[line width=.7pt,rounded corners=2pt] (2.57,5.29) rectangle (3.26,5.98);
\draw[line width=1.5pt] (2.6252,5.29) -- (2.6803999999999997,5.29);
\draw[line width=.95pt,line join=round] (3.0667999999999997,5.386393)--(3.159398,5.554753)--(2.9742019999999996,5.554753)--cycle;
\draw[line width=.7pt,rounded corners=2pt] (3.46,5.29) rectangle (4.15,5.98);
\draw[line width=1.5pt] (3.5152,5.29) -- (3.5704,5.29);
\draw[line width=.95pt] (3.6532,5.4832) circle[radius=0.08417999999999999in];
\draw[line width=.95pt,line join=round] (3.9568,5.386393)--(4.049398,5.554753)--(3.8642019999999997,5.554753)--cycle;
\draw[line width=.95pt] (3.56902,5.7026200000000005) rectangle (3.73738,5.87098);
\draw[line width=.7pt,rounded corners=2pt] (4.35,5.29) rectangle (5.039999999999999,5.98);
\draw[line width=1.5pt] (4.4052,5.29) -- (4.4604,5.29);
\draw[line width=.95pt] (4.5432,5.4832) circle[radius=0.08417999999999999in];
\draw[line width=.7pt,rounded corners=2pt] (5.239999999999999,5.29) rectangle (5.93,5.98);
\draw[line width=1.5pt] (5.2951999999999995,5.29) -- (5.3504,5.29);
\draw[line width=.95pt,line join=round] (5.7368,5.386393)--(5.829397999999999,5.554753)--(5.644202,5.554753)--cycle;
\draw[line width=.95pt] (5.349019999999999,5.7026200000000005) rectangle (5.517379999999999,5.87098);
\draw[line width=.7pt,rounded corners=2pt] (6.13,5.29) rectangle (6.82,5.98);
\draw[line width=1.5pt] (6.1852,5.29) -- (6.2404,5.29);
\draw[line width=.95pt] (6.3232,5.4832) circle[radius=0.08417999999999999in];
\draw[line width=.95pt,line join=round] (6.6268,5.386393)--(6.719398,5.554753)--(6.5342020000000005,5.554753)--cycle;
\draw[line width=.7pt,rounded corners=2pt] (7.02,5.29) rectangle (7.709999999999999,5.98);
\draw[line width=1.5pt] (7.0752,5.29) -- (7.1304,5.29);
\draw[line width=.95pt] (7.12902,5.7026200000000005) rectangle (7.2973799999999995,5.87098);
\end{scope}\end{tikzpicture}\mbox{}
\newpage
\begin{tikzpicture}[remember picture,overlay]\begin{scope}[shift={(current page.north west)},x=1in,y=-1in]
\node[anchor=north west,inner sep=0pt,text width=7.2in,align=left,font=\fontsize{11}{13.75}\selectfont] at (0.65,0.41) {Week 19 / No card inside another / Grades 4--5};
\draw[line width=.35pt] (.65,.70)--(7.85,.70);
\node[anchor=north west,inner sep=0pt,text width=6.7in,align=left,font=\fontsize{9}{11.25}\selectfont] at (0.65,10.43) {Bellingham Math Circle / Week 19 / F19-45-v2};
\node[anchor=north west,inner sep=0pt,text width=0.55in,align=right,font=\fontsize{9}{11.25}\selectfont] at (7.3,10.43) {5};
\node[anchor=north west,inner sep=0pt,text width=7.2in,align=left,font=\fontsize{14}{17.5}\selectfont] at (0.65,0.98) {\textbf{Problem 5:} Arrange all sixteen cards from Problem 2 into rows, with each card fitting inside the next card in its row. Use every card exactly once and make as few rows as possible.};
\draw[black!22,line width=.35pt] (.65,7.13)--(7.85,7.13);
\node[anchor=north west,inner sep=0pt,text width=7.2in,align=left,font=\fontsize{14}{17.5}\selectfont] at (0.65,7.43) {\textbf{Problem 6:} What is the largest possible collection from the sixteen-card deck? Give a collection of that size and explain why every larger collection breaks the rule.};
\end{scope}\end{tikzpicture}\mbox{}
\newpage
\begin{tikzpicture}[remember picture,overlay]\begin{scope}[shift={(current page.north west)},x=1in,y=-1in]
\node[anchor=north west,inner sep=0pt,text width=7.2in,align=left,font=\fontsize{11}{13.75}\selectfont] at (0.65,0.41) {Week 19 / No card inside another / Grades 4--5};
\draw[line width=.35pt] (.65,.70)--(7.85,.70);
\node[anchor=north west,inner sep=0pt,text width=6.7in,align=left,font=\fontsize{9}{11.25}\selectfont] at (0.65,10.43) {Bellingham Math Circle / Week 19 / F19-45-v2};
\node[anchor=north west,inner sep=0pt,text width=0.55in,align=right,font=\fontsize{9}{11.25}\selectfont] at (7.3,10.43) {6};
\node[anchor=north west,inner sep=0pt,text width=7.2in,align=left,font=\fontsize{14}{17.5}\selectfont] at (0.65,0.98) {\textbf{Problem 7:} Use the sixteen-card deck. Find the largest collection that includes the card shown. Explain why no larger collection can include it.};
\draw[line width=.7pt,rounded corners=2pt] (3.765,2.0) rectangle (4.735,2.9699999999999998);
\draw[line width=1.5pt] (3.8426,2.0) -- (3.9202000000000004,2.0);
\draw[line width=.95pt] (4.0366,2.2716) circle[radius=0.11834in];
\draw[black!22,line width=.35pt] (.65,5.49)--(7.85,5.49);
\node[anchor=north west,inner sep=0pt,text width=7.2in,align=left,font=\fontsize{14}{17.5}\selectfont] at (0.65,5.79) {\textbf{Problem 8:} Find every collection that reaches your largest size from Problem 6. Explain why your list is complete.};
\end{scope}\end{tikzpicture}\mbox{}
\end{document}
